% !TeX root = ../main.tex
% ============================================================
%  appendix/app_proofs.tex
% ============================================================

\section{Proofs}
\label{app:proofs}

%-----------------------------------------------------------------
\begin{proof}[Proof of \Cref{lem:lb_validity}]
Fix $k\ge0$ and let $(\taxindex^{*},\bz^{*})$ be optimal
for~\eqref{eq:leader}, so $\bz^{*}\in\Sol{\taxindex^{*}}$. The feasible
region $\setX$ of the follower's problem does not depend on $\taxindex$
(\Cref{sec:model:follower}); only the objective~\eqref{eq:follower}
does. Hence every $\bz^{(j)}\in\Sol{\taxindex^{(j)}}$ used to generate
cut $j$, for $j$ ranging over the cut set $\mathcal{C}$ accumulated by
iteration $k$, is itself a feasible -- if possibly suboptimal -- point
of the follower's problem at $\taxindex^{*}$. Optimality of $\bz^{*}$ at
$\taxindex^{*}$ therefore gives
\begin{equation*}
\Cost(\bz^{*})+\gridstep\,\taxindex^{*}\,\Emis(\bz^{*})
\;\le\;
\Cost\!\bigl(\bz^{(j)}\bigr)+\gridstep\,\taxindex^{*}\,\Emis\!\bigl(\bz^{(j)}\bigr)
= c_{j}+\gridstep\,\taxindex^{*}\,e_{j}
\qquad\text{for every }(\taxindex^{(j)},c_{j},e_{j})\in\mathcal{C}.
\end{equation*}
Thus $(\taxindex^{*},\bz^{*})$ satisfies cut~\eqref{eq:cut} for every
cut accumulated by iteration $k$, i.e.\ it is feasible for the
relaxation master defining $\LB^{(k)}$ in~\eqref{eq:lb_def}. Since
$\LB^{(k)}$ minimises $\Emis(\cdot)$ over that (larger) feasible set,
$\LB^{(k)}\le\Emis(\bz^{*})$.
\end{proof}

%-----------------------------------------------------------------
\begin{proof}[Proof of \Cref{lem:tiebreak_validity}]
$\bz'\in\Sol{\taxindex'}$ is, in particular, a feasible point of
$\setX$, and $\setX$ does not depend on $\taxindex$
(\Cref{sec:model:follower}). Fix any cut $(\taxindex^{(j)},c_{j},e_{j})
\in\mathcal{C}$ accumulated by iteration $k$. As in the proof of
\Cref{lem:lb_validity}, $\bz^{(j)}$ is itself a feasible -- if
possibly suboptimal -- point of the follower's problem at $\taxindex'$,
so optimality of $\bz'$ at $\taxindex'$ gives
\begin{equation*}
\Cost(\bz')+\gridstep\,\taxindex'\,\Emis(\bz')
\;\le\;
\Cost\!\bigl(\bz^{(j)}\bigr)+\gridstep\,\taxindex'\,\Emis\!\bigl(\bz^{(j)}\bigr)
= c_{j}+\gridstep\,\taxindex'\,e_{j}.
\end{equation*}
Since this holds for every cut in $\mathcal{C}$, the pair
$(\taxindex',\bz')$ satisfies~\eqref{eq:cut} for every accumulated cut,
i.e.\ it is feasible for the Stage-1 master defining $\LB^{(k)}$.
Together with $\Emis(\bz')=\LB^{(k)}$ (hypothesis), $(\taxindex',\bz')$
is therefore one of the master points Stage~2 minimises $\taxindex$
over. Since $\taxindex^{\star}$ is the minimum of $\taxindex$ over that
same feasible set, and $\taxindex'$ is feasible for it,
$\taxindex^{\star}\le\taxindex'$.
\end{proof}

%-----------------------------------------------------------------
\begin{proof}[Proof of \Cref{thm:convergence}]
\textbf{Validity of bounds.}
By \Cref{lem:lb_validity,lem:ub_validity}, the inequalities
$\LB^{(k)}\le\Emis(\bz^{*})\le\UB^{(k)}$ hold for every iteration $k$
of \Cref{alg:dbp}.

\textbf{Finite termination.}
Fix any tax index $\taxindex'\in\mathcal{T}^{\tau}$ already evaluated by
iteration $k$, i.e.\ $\taxindex'=\taxindex^{(j)}$ for some cut
$(\taxindex^{(j)},c_{j},e_{j})\in\mathcal{C}$. Evaluating
cut~\eqref{eq:cut} at $\taxindex=\taxindex'$ gives
$\Cost(\bz)+\gridstep\,\taxindex'\,\Emis(\bz)\le
c_{j}+\gridstep\,\taxindex'\,e_{j}$ for every $\bz$ feasible in the
relaxation master defining $\LB^{(k)}$; and since
$c_{j}+\gridstep\,\taxindex'\,e_{j}=\Cost(\bz^{(j)})+\gridstep\,\taxindex'\,\Emis(\bz^{(j)})$
is itself the optimal value of~\eqref{eq:follower} at $\taxindex'$
(because $\bz^{(j)}\in\Sol{\taxindex'}$), no feasible $\bz$ can improve
on $\bz^{(j)}$ in objective~\eqref{eq:follower} at $\taxindex'$. Hence
any $\bz$ chosen by the master at $\taxindex=\taxindex'$ must itself lie
in $\Sol{\taxindex'}$, and minimising $\Emis(\cdot)$ over that
restriction recovers exactly $\min\{\Emis(\bz):\bz\in\Sol{\taxindex'}\}$,
the optimistic value at $\taxindex'$ (\Cref{as:opt}) -- i.e.\ the
master's bound at every previously-evaluated index is \emph{exact},
not merely a relaxation. Since $\mathcal{T}^{\tau}=\{\taxindexlb,\ldots,\taxindexub\}$
has only $\taxindexub-\taxindexlb+1$ elements, after at most $\taxindexub-\taxindexlb+1$
iterations every admissible tax index either (i) has been evaluated
exactly, making $\LB^{(k)}$ exact there, or (ii) is dominated, in the
master's emissions-minimisation, by an index covered by (i). In
either case $\LB^{(k)}=\UB^{(k)}$ once $k=\taxindexub-\taxindexlb+1$, so the gap
closes within a finite, explicitly bounded number of iterations --
with no continuity or breakpoint-density argument on
$\Emis(\bz^{*}(\taxindex))$ required.

\textbf{$\epsilon$-optimality.}
On termination, $\gap=(\UB^{(k)}-\LB^{(k)})/\UB^{(k)}\le\epsilon$, so
the returned index $\taxindex^{\dagger}$ satisfies
$\Emis\!\bigl(\bz^{\dagger}\bigr)=\UB^{(k)}
 \le\LB^{(k)}/(1-\epsilon)
 \le\Emis(\bz^{*})/(1-\epsilon)$, establishing the claimed
$\epsilon$-optimality guarantee.
\end{proof}
